Activation-space interventions are increasingly used to control how varied a language model's
outputs are, on the premise that output diversity is governed by a direction that can be
estimated by contrasting high- and low-dispersion generations and then added back at inference
time. We test that premise directly. Across seven open-weight models spanning four families
(410M to 7B parameters), we find that a direction fitted this way has no stable causal role:
injecting it across 35 dose levels never increases diversity, the median within-model
dose--response slope is negative (Spearman $\rho = -0.30$, $p = 0.016$, 0/7 models positive),
and at matched coefficients it is indistinguishable from---numerically worse than---a random
vector ($-0.082$ \vendi{}, $p = 0.297$). The reason is that the estimator does not measure what
it is assumed to measure. A split-half test shows that the pooled direction is reproducible only
because it encodes prompt identity: for pythia-410m the split-half cosine is $+0.999$ when
prompts are mixed but $-0.037$ when the two halves come from the same prompt, and this pattern
holds for all seven models, where within-prompt agreement never exceeds $|0.12|$. The pooled
vector is therefore a sum of near-orthogonal prompt-local components, none of which predicts the
others. Ablating the response subspace still collapses diversity, but destroying the pairing that
defines the contrast does the same, so the effect reflects removal of generative capacity rather
than control of a diversity axis. We conclude that output diversity has no linearly decodable
direction to steer, and we describe the placebo and quality controls a claim of successful
diversity steering would need to pass.
